% Source: https://artofproblemsolving.com/community/c3h1644250_mirror_problem
% Stable identity: aops-topic-1644250-q01
For a stack of $k$ consecutive mirrors, let $T_k$ be the probability that a photon arriving at the stack from one side eventually leaves it on the other side (is transmitted), and $R_k$ the probability that it eventually comes back out on the side it came from. The stack looks the same from both sides, so these do not depend on the side. A photon inside a finite stack leaves it with probability $1$ (at every hit it has probability at least $\dfrac{1}{2}$ of making progress towards an exit, and a finite chain cannot avoid the exits forever), so $R_k=1-T_k$. Also $T_1=\dfrac{1}{2}$.

\textbf{Adding a mirror.} Let a photon arrive from the left at the stack of $k+1$ mirrors. To get through it must first pass the first $k$ mirrors (probability $T_k$). Then it hits mirror $k+1$: it passes (probability $\dfrac{1}{2}$), or it is reflected back into the first $k$ mirrors (probability $\dfrac{1}{2}$), from which it returns towards mirror $k+1$ with probability $R_k$ (otherwise it escapes to the left). Summing over the number $j$ of such round trips,
\[
T_{k+1}=T_k\cdot\dfrac{1}{2}\sum_{j\ge0}\left(\dfrac{R_k}{2}\right)^j=\dfrac{T_k}{2-R_k}=\dfrac{T_k}{1+T_k}.
\]
Taking reciprocals, $\dfrac{1}{T_{k+1}}=\dfrac{1}{T_k}+1$, and since $\dfrac{1}{T_1}=2$, induction gives $\dfrac{1}{T_k}=k+1$, that is $T_k=\dfrac{1}{k+1}$.

\textbf{The wall.} The photon starts moving right towards the whole stack. If it is transmitted it reaches the detector. If it ever comes back out to the left of the first mirror, it hits the wall and is destroyed. So the probability of detection is
\[
T_{2047}=\boxed{\dfrac{1}{2048}}.
\]
